% Part: counterfactuals
% Chapter: introduction
% Section: strict-conditional

\documentclass[../../../include/open-logic-section]{subfiles}

\begin{document}

\olfileid{cnt}{int}{str}

\olsection{కఠిన సోపాధికం}

లూయిస్ \emph{కఠిన సోపాధికం}~$\strictif$ను
ప్రవేశపెట్టి, భౌతిక సోపాధికం కాకుండా ఇదే
పర్యవసాన భావానికి సరిపోతుందని వాదించారు.
అలేథిక్ మోడల్ తర్కంలో $!A \strictif !B$ను
$\Box(!A \lif !B)$గా నిర్వచించవచ్చు.
కాబట్టి సంబంధిత భౌతిక సోపాధికం
అనివార్యమైనప్పుడే కఠిన సోపాధికం
(ఒక లోకంలో) సత్యం.

భౌతిక సోపాధికం వైరుధ్యాభాసాల విషయంలో
కఠిన సోపాధికం ఎలా పనిచేస్తుంది?
అసత్య పూర్వపక్షం గల కఠిన సోపాధికం,
సత్య ఉత్తరపక్షం గల కఠిన సోపాధికం
సత్యం కావచ్చు, అసత్యమూ కావచ్చు.
అంతేకాకుండా $(!A \strictif !B) \lor
(!B \strictif !A)$ సర్వత్ర చెల్లదు.
కఠిన సోపాధికం $!A \strictif !B$,
$\lnot !A \lor !B$కు తుల్యమూ కాదు;
అందువల్ల అది సత్యమూల్యాధారితం కాదు.

మనకు ఇవి ఉన్నాయి:
\begin{align}
  !A \strictif !B & \Entails \lnot !A \lor !B
  \text{ కానీ:}\\
  \lnot !A \lor !B & \Entails/ !A \strictif !B\\
  !B & \Entails/ !A \strictif !B\\
  \lnot !A & \Entails/ !A \strictif !B\\
  \lnot(!A \strictif !B) & \Entails/ !A \land \lnot !B
  \text{ కానీ:}\\
  !A \land \lnot !B & \Entails \lnot(!A \strictif !B)
  \intertext{అయినా కఠిన సోపాధికానికి
    మోడస్ పోనెన్స్ వర్తిస్తుంది:}
  !A, !A \strictif !B & \Entails !B
  \intertext{కఠిన సోపాధికాలను
    ఒకటిగా చేర్చవచ్చు:}
  !A \strictif !B,  !A \strictif !C & \Entails !A \strictif (!B \land !C)
  \intertext{కఠిన సోపాధికానికి
    పూర్వపక్ష బలపరచడం వర్తిస్తుంది:}
  !A \strictif !B & \Entails (!A \land !C) \strictif !B
  \intertext{కఠిన సోపాధికం కూడా సంక్రామకం:}
  !A \strictif !B, !B \strictif !C & \Entails !A \strictif !C
  \intertext{చివరగా, కఠిన సోపాధికం తన
    ప్రతివిపర్యయానికి తుల్యం:}
  !A \strictif !B & \Entails \lnot !B \strictif \lnot !A\\
  \lnot !B \strictif \lnot !A & \Entails !A \strictif !B
\end{align}
\sourcecorrection{OLTECNTSTR-001}{మూల ప్రదర్శనలో ఐదో అపర్యవసానానికి భౌతిక సోపాధిక నిరాకరణను వాడారు; అది నిజానికి పూర్వపక్షం-నిరాకృత ఉత్తరపక్షం సంయోగానికి తుల్యం. ఇక్కడ ఉద్దేశించిన అపర్యవసానం కఠిన సోపాధిక నిరాకరణది; దాని సంకేతాన్ని సరిచేశాం.}

\begin{prob}
  కఠిన సోపాధికానికి వర్తించని
  పర్యవసాన సంబంధాలకు
  \Log{S5}-ప్రతిదృష్టాంతాలు ఇవ్వండి; అంటే:
  \begin{enumerate}
  \item $\lnot p \Entails/ \Box(p \lif q)$
  \item $q \Entails/ \Box(p \lif q)$
  \item $\lnot\Box(p \lif q) \Entails/ p \land \lnot q$
  \item $\Entails/ \Box(p \lif q) \lor \Box(q \lif p)$
  \end{enumerate}
\end{prob}

\begin{prob}
  కఠిన సోపాధికానికి వర్తించే
  పర్యవసాన సంబంధాలను కింది వాటికి
  \Log{S5}-నిరూపణలు ఇచ్చి చూపండి:
  \begin{enumerate}
  \item  $\Box(!A \lif !B) \Entails \lnot !A \lor !B$
  \item $!A \land \lnot !B \Entails \lnot\Box(!A \lif !B)$
  \item $!A, \Box(!A \lif !B) \Entails !B$
  \item $\Box(!A \lif !B), \Box(!A \lif !C) \Entails \Box(!A \lif (!B
    \land !C))$
  \item $\Box(!A \lif !B) \Entails \Box((!A \land !C) \lif !B)$
  \item $\Box(!A \lif !B), \Box(!B \lif !C) \Entails \Box(!A
    \lif !C)$
  \item $\Box(!A \lif !B) \Entails \Box(\lnot !B \lif \lnot !A)$
  \item $\Box(\lnot !B \lif \lnot !A) \Entails \Box(!A \lif !B)$
  \end{enumerate}
  \end{prob}

అయితే కఠిన సోపాధికానికి కూడా దాని
సొంత ``వైరుధ్యాభాసాలు'' ఉన్నాయి.
అసత్య పూర్వపక్షం లేదా సత్య ఉత్తరపక్షం
గల భౌతిక సోపాధికం సత్యమైనట్టే,
\emph{అనివార్యంగా} అసత్య పూర్వపక్షం
లేదా అనివార్యంగా సత్య ఉత్తరపక్షం గల
కఠిన సోపాధికం సత్యం. అంతేకాకుండా,
సత్యమైన ఏ కఠిన సోపాధికమైనా అనివార్యంగా
సత్యం; అసత్యమైన ఏ కఠిన సోపాధికమైనా
అనివార్యంగా అసత్యం. మరో మాటలో:
\begin{align}
  \Box \lnot !A & \Entails !A \strictif !B\\
  \Box !B & \Entails !A \strictif !B\\
  !A \strictif !B& \Entails \Box(!A \strictif !B)\\
  \lnot(!A \strictif !B) & \Entails \Box\lnot(!A \strictif !B)
\end{align}
మీరు $\strictif$ను ``పర్యవసిస్తుంది''
అని భావిస్తే ఇవి సమస్యలు కావు. తార్కిక
పర్యవసాన సంబంధాలు గణిత వాస్తవాలే;
అందువల్ల అవి యాదృచ్ఛిక స్థితులపై
ఆధారపడవు. కానీ ``ఒకవేళ \dots అయితే
\dots'' భావాన్ని పట్టుకోవాల్సిన తార్కిక
సంయోజకంగా $\strictif$ను వాడాలంటే,
ముఖ్యంగా చివరి రెండు వాదనలు సమస్యలు
రేకెత్తిస్తాయి. ఎందుకంటే అనివార్యం
కాకుండానే సత్యమైన లేదా అసత్యమైన
``ఒకవేళ \dots అయితే \dots'' వాక్యాలు
నిశ్చయంగా ఉన్నాయి; నిజానికి అవి
సాధారణంగా అనివార్యాలూ కావు,
అసంభవాలూ కావు.

\begin{prob}
  \Log{S5}లో కింది వాటికి నిరూపణలు ఇవ్వండి:
  \begin{enumerate}
    \item  $\Box \lnot !A \Entails !A \strictif !B$
    \item $!A \strictif !B \Entails \Box(!A \strictif !B)$
    \item $\lnot(!A \strictif !B)  \Entails \Box\lnot(!A \strictif !B)$
  \end{enumerate}
  అందుకు $\strictif$ నిర్వచనాన్ని వాడండి.
\end{prob}

\end{document}
